\clearpage\pagestyle{otMain}\fontsize{11}{15}\selectfont\corrigesChapitre{12}{Pythagore et distances}
\begin{blocLivre}{cor-12-1}\begin{corrige}{1}{Reconnaître l’hypoténuse}\demarche{Repérer d’abord le sommet de l’angle droit.}
Le côté opposé à A est $\res{BC}$. L’égalité est $\res{BC^2=AB^2+AC^2}$. Le nom du sommet rectangle détermine le bon côté.
\label{fig-cor-12-1}\figureGeo[corrige]{pyth-01}
\end{corrige}\end{blocLivre}
\begin{blocLivre}{cor-12-2}\begin{corrige}{2}{Un triangle classique}\demarche{Additionner les carrés pour trouver l’hypoténuse.}
Pythagore donne $BC^2=6^2+8^2=100$. Comme BC est positive, $\res{BC=10\text{ cm}}$.
\label{fig-cor-12-2}\figureGeo[corrige]{pyth-02}
\end{corrige}\end{blocLivre}
\begin{blocLivre}{cor-12-3}\begin{corrige}{3}{Le côté manquant}\demarche{Isoler le carré inconnu avant d’extraire la racine.}
$AC^2=13^2-5^2=144$, donc $\res{AC=12\text{ cm}}$. AC est inférieure à l’hypoténuse 13 cm.
\label{fig-cor-12-3}\figureGeo[corrige]{pyth-03}
\end{corrige}\end{blocLivre}
\begin{blocLivre}{cor-12-4}\begin{corrige}{4}{Garder une racine}\demarche{Arrondir uniquement le résultat final.}
$BC^2=16+49=65$, donc $\res{BC=\sqrt{65}\text{ cm}}\approx\res{8{,}1\text{ cm}}$. L’arrondi n’est pas une égalité exacte.
\label{fig-cor-12-4}\figureGeo[corrige]{pyth-04}
\end{corrige}\end{blocLivre}
\begin{blocLivre}{cor-12-5}\begin{corrige}{5}{Tester la réciproque}\demarche{Comparer les deux membres séparément.}
BC est le plus grand côté. $15^2=225$ et $9^2+12^2=225$. Par la réciproque, ABC est $\res{\text{rectangle en A}}$.
\label{fig-cor-12-5}\figureGeo[corrige]{pyth-05}
\end{corrige}\end{blocLivre}
\begin{blocLivre}{cor-12-6}\begin{corrige}{6}{Refuser une fausse égalité}\demarche{Une égalité approchée ne suffit pas.}
$BC^2=81$ alors que $AB^2+AC^2=25+49=74$. Ces nombres diffèrent ; par la contraposée, $\res{\text{ABC n’est pas rectangle}}$.
\label{fig-cor-12-6}\figureGeo[corrige]{pyth-06}
\end{corrige}\end{blocLivre}
\begin{blocLivre}{cor-12-7}\begin{corrige}{7}{Une échelle contre un mur}\demarche{Justifier l’angle droit par le mur et le sol.}
ABC est rectangle en A. $AC^2=5^2-1{,}4^2=23{,}04$, donc $\res{AC=4{,}8\text{ m}}$. La hauteur reste inférieure à la longueur de l’échelle.
\label{fig-cor-12-7}\figureGeo[corrige]{pyth-07}
\end{corrige}\end{blocLivre}
\begin{blocLivre}{cor-12-8}\begin{corrige}{8}{Diagonale d’un rectangle}\demarche{Séparer longueur intérieure et contour.}
$BC^2=15^2+20^2=625$, donc la diagonale vaut $\res{25\text{ cm}}$. Le périmètre du rectangle est $2(15+20)=\res{70\text{ cm}}$. Il ne faut pas y ajouter la diagonale intérieure.
\label{fig-cor-12-8}\figureGeo[corrige]{pyth-08}
\end{corrige}\end{blocLivre}
\begin{blocLivre}{cor-12-9}\begin{corrige}{9}{Deux points du plan}\demarche{Construire les deux déplacements perpendiculaires.}
Les écarts horizontal et vertical valent 6 et 8. Dans le triangle rectangle construit, $AB^2=6^2+8^2=100$, d’où $\res{AB=10\text{ cm}}$.
\label{fig-cor-12-9}\figureGeo[corrige]{pyth-repere}
\end{corrige}\end{blocLivre}
\begin{blocLivre}{cor-12-10}\begin{corrige}{10}{Un angle à identifier}\demarche{Nommer le sommet opposé à l’hypoténuse.}
Le plus grand côté est AB. $AB^2=100=36+64=AC^2+BC^2$. La réciproque donne un angle droit $\res{\text{en C}}$, opposé à AB. La position du dessin ne détermine pas l’angle droit.
\label{fig-cor-12-10}\figureGeo[corrige]{pyth-10}
\end{corrige}\end{blocLivre}
\begin{blocLivre}{cor-12-11}\begin{corrige}{11}{Une voile et son aire}\demarche{Enchaîner une longueur puis une aire.}
$AC^2=16-5{,}76=10{,}24$, donc AC = 3,2 m. L’aire vaut $AB\times AC/2=2{,}4\times3{,}2/2=\res{3{,}84\text{ m}^2}$. On utilise les deux côtés perpendiculaires.
\label{fig-cor-12-11}\figureGeo[corrige]{pyth-11}
\end{corrige}\end{blocLivre}
\begin{blocLivre}{cor-12-12}\begin{corrige}{12}{Passer un panneau}\demarche{Comparer des longueurs dans le même modèle.}
La diagonale vaut $\sqrt{0{,}9^2+1{,}2^2}=\res{1{,}5\text{ m}}$. Comme $1{,}48<1{,}5$, la tige peut être placée suivant cette diagonale dans le modèle proposé. L’épaisseur réelle pourrait imposer une marge.
\label{fig-cor-12-12}\figureGeo[corrige]{pyth-12}
\end{corrige}\end{blocLivre}
\begin{blocLivre}{cor-12-13}\begin{corrige}{13}{Une rampe}\demarche{Distinguer avancée horizontale et longueur inclinée.}
$BC=\sqrt{10^2+0{,}75^2}\approx\res{10{,}03\text{ m}}$. La pente vaut $100\times0{,}75/10=\res{7{,}5\,\%}$. On ne divise pas la hauteur par la longueur inclinée dans cette définition.
\label{fig-cor-12-13}\figureGeo[corrige]{pyth-13}
\end{corrige}\end{blocLivre}
\begin{blocLivre}{cor-12-14}\begin{corrige}{14}{Deux chemins}\demarche{Comparer les trajets après les avoir calculés.}
Le trajet par les côtés vaut 280 m. La diagonale vaut $\sqrt{120^2+160^2}=200$ m. L’économie est $\res{80\text{ m}}$, soit $80/280=2/7$ du trajet initial, environ 28,6 \%.
\label{fig-cor-12-14}\figureGeo[corrige]{pyth-parc}
\end{corrige}\end{blocLivre}
\begin{blocLivre}{cor-12-15}\begin{corrige}{15}{Une hauteur de triangle}\demarche{Justifier la perpendicularité avant Pythagore.}
Dans un triangle isocèle, la médiane CH est aussi une hauteur. AH = 5 cm et ACH est rectangle en H. $CH^2=13^2-5^2=144$, donc CH = 12 cm. L’aire vaut $10\times12/2=\res{60\text{ cm}^2}$.
\label{fig-cor-12-15}\figureGeo[corrige]{pyth-iso}
\end{corrige}\end{blocLivre}
